\chapter{Testing Speaker Diarisation Results}
\label{app:speaker_diarisation_results}
\subsection{Doctor Transcript}
\noindent
{[Doctor 0.79]} This research investigates the application of on-device machine learning and clinical administration \newline
{[Doctor 0.70]} through the development of an offline Windows application. As a proof of concept, \newline
{[Doctor 0.84]} the primary objective is to evaluate the feasibility of a native Windows application \newline
{[Doctor 0.74]} running on an Intel AI PC to autonomously record consultations, transcribe audio, \newline
{[Doctor 0.78]} perform speaker diarisation, generate clinical summaries, and retrieve medical guidance within an entirely offline environment. \newline
{[Doctor 0.75]} While existing cloud-based systems such as TORTUS AI, Heidi Health, and Nuance DAX Copilot \newline
{[Doctor 0.72]} highlight the utility of automated clinical documentation within healthcare, \newline
{[Doctor 0.78]} they rely on networked deployment, extensive integration, and external data processing. \newline
{[Doctor 0.76]} In the context of the National Health Service, these dependencies introduce significant data governance and infrastructural challenges. \newline
{[Doctor 0.74]} Consequently, developing an offline proof of concept provides a mechanism to evaluate an alternative deployment architecture designed to address these specific operational constraints.

\vspace{1em}

\subsection*{Patient Transcript}
\noindent
{[Patient 0.42]} This research investigates the application of on-device machine learning \newline
{[Patient 0.63]} and clinical administration through the development of an offline Windows application. \newline
{[Patient 0.51]} As a proof of concept, the primary objective is to evaluate the feasibility \newline
{[Patient 0.58]} of a native Windows application running on an Intel AI PC to autonomously \newline
{[Patient 0.47]} record consultations, transcribe audio, perform speaker diarisation, generate clinical summaries, \newline
{[Patient 0.55]} and retrieve medical guidance within an entirely offline environment. While existing cloud-based \newline
{[Patient 0.63]} systems such as TORTUS AI, Heidi Health, and Nuance DAX Copilot highlight the utility \newline
{[Patient 0.49]} of automated clinical documentation within healthcare, they rely on networked deployment, extensive integration, and external data processing. \newline
{[Patient 0.60]} In the context of the National Health Service, these dependencies introduce significant data governance \newline
{[Patient 0.52]} and infrastructural challenges. Consequently, developing an offline proof of concept provides a mechanism to evaluate an alternative deployment architecture designed to address these specific operational constraints.